\documentclass[11pt]{article}
\usepackage[margin=1in]{geometry}
\usepackage{amsmath,amssymb,mathtools}
\usepackage{fontspec,unicode-math}
\setmainfont{Latin Modern Roman}
\setsansfont{Latin Modern Sans}
\setmathfont{Latin Modern Math}
\usepackage{xurl,hyperref}
\hypersetup{colorlinks=true,urlcolor=blue,linkcolor=blue}
\newcommand{\sref}[2]{\hyperlink{source:#1:#2}{[S#2]}}
\newcommand{\cataloguescope}[1]{\paragraph{Recorded mathematical question.}#1}
\begin{document}
\section{Resistance growth on the UIPT}\label{statement}

\subsection{Definitions and mathematical statement}

\paragraph{Shared notation.}$\mathbb N=\{1,2,\ldots\}$, and $\mathbb Z,\mathbb Q,\mathbb R,\mathbb C$ denote the integers, rationals, reals and complex numbers. Any other conventions are specified locally.


The uniform infinite planar triangulation (UIPT) is the random rooted locally finite infinite triangulation obtained as the distributional local limit of uniform finite rooted triangulations as their sizes tend to infinity. Give every edge resistance $1$. For its root $o$, let $R(r)$ be effective resistance from $o$ to the vertices outside the graph-distance ball $B_G(o,r-1)$, with those exterior vertices wired into one terminal. Equivalently this is resistance to the radius-$r$ boundary with that boundary wired.
\paragraph{Statement recorded in the supplied source.}
Determine the almost-sure asymptotic growth of $R(r)$ as $r\to\infty$, including its order of magnitude and any asymptotic exponent or leading expression. \sref{statement}{2}
\subsection{Short English statement}
At what rate does electrical resistance to a distant boundary grow in the UIPT?
\subsection{Sources}
\begin{description}
\item[\hypertarget{source:statement:1}{S1}] ULAM.AI and the UnsolvedMath Contributors, \emph{UnsolvedMath: A Curated Collection of Open Mathematics Problems}, catalogue entry. CC BY 4.0. \url{https://huggingface.co/datasets/ulamai/UnsolvedMath}.
\item[\hypertarget{source:statement:2}{S2}] Itai Benjamini, \emph{Random Planar Metrics}; original author manuscript preserved in Arquivo.pt capture of 31 December 2020, Section 3, item (1), p. 3. \url{https://arquivo.pt/noFrame/replay/20201231041544id_/http://www.wisdom.weizmann.ac.il/~itai/randomplanar2.pdf}.
\end{description}
\label{end:statement}
\end{document}
